\documentclass{ibetest}
\begin{document}
\maketest{Progress Test 3.A --- Elec-S1}
         {3.A \,$\cdot$\, Kirchhoff's laws (combinations, divider, Millman)}
         {10 minutes}{14}

% ---------------------------------------------------------------------------
\exercise{Equivalent resistance}{1 mark}
Give the equivalent resistance of $R_1 = \qty{6.0}{\kilo\ohm}$ and
$R_2 = \qty{3.0}{\kilo\ohm}$ connected in parallel.
\answer{1}{2.2cm}{%
  $\Req = \dfrac{R_1 R_2}{R_1 + R_2} = \dfrac{6.0\times3.0}{6.0 + 3.0}\ \mathrm{k\Omega}
   = \dfrac{18}{9.0}\ \mathrm{k\Omega}$ \qquad $\boxed{\Req = 2.0\ \mathrm{k\Omega}}$
  \quad (smaller than the smaller one)}
\begin{scheme}
\pt{1}{1 pt}{$\Req = 2.0\ \mathrm{k\Omega}$ (``product over sum'').}
\end{scheme}

% ---------------------------------------------------------------------------
\exercise{Kirchhoff's laws}{5 marks}
State Kirchhoff's node law and mesh law. For each law, give the fundamental principle it
expresses.
\answer{2,2,1}{6.2cm}{%
  \textbf{Node law.} At any node, the sum of the currents entering the node equals the sum
  of the currents leaving it:
  $\sum i_{\mathrm{in}} = \sum i_{\mathrm{out}}$, i.e.\ $\sum_k \varepsilon_k i_k = 0$
  with $\varepsilon_k = +1$ for a current entering and $-1$ for a current leaving.\\[4pt]
  \textbf{Mesh law.} Along any mesh (closed path), the algebraic sum of the voltages is
  zero: $\sum_k \varepsilon_k u_k = 0$, with $\varepsilon_k = +1$ if the arrow of $u_k$
  points in the direction of travel and $-1$ otherwise.\\[4pt]
  The node law expresses the \textbf{conservation of electric charge} (no charge piles up at
  a node). The mesh law expresses the \textbf{conservation of energy} (after a closed path the
  potential returns to its starting value).}
\begin{scheme}
\pt{1}{2 pts}{Node law: statement (1~pt) and algebraic form with signs (1~pt).}
\pt{2}{2 pts}{Mesh law: statement (1~pt) and sign convention along the path (1~pt).}
\pt{3}{1 pt}{Charge conservation \emph{and} energy conservation, each matched to its law.}
\end{scheme}
\begin{tip}
Box~6 of the handout is easy to misread: its brackets refer to the \emph{first} sentence of
each paragraph. The same current through components in series comes from the node law; the
same voltage across components in parallel comes from the mesh law. Currents add up at a
node (node law); voltages add up round a loop (mesh law).
\end{tip}

% ---------------------------------------------------------------------------
\exercise{Quick circuit laws}{3 marks}
\question{$R_1 = \qty{1.0}{\kilo\ohm}$ and $R_2 = \qty{3.0}{\kilo\ohm}$ are in series under
  $u = \qty{12}{V}$. The voltage $u_2$ across $R_2$ is:}
\opts{0}{\qty{3.0}{V}}{0}{\qty{4.0}{V}}{1}{\qty{9.0}{V}}{0}{\qty{12}{V}}
\why{voltage divider: $u_2 = \frac{R_2}{R_1 + R_2}\,u = \frac{3}{4}\times12 = \qty{9.0}{V}$.
  The larger resistance takes the larger share.}
\question{A circuit has $N = 3$ electrically distinct nodes and $B = 5$ branches. The number of
  independent mesh equations is:}
\opts{0}{2}{1}{3}{0}{5}{0}{8}
\why{$m = B - N + 1 = 5 - 3 + 1 = 3$ (and there are $N - 1 = 2$ independent node
  equations: $2 + 3 = 5$ equations for the 5 unknown branch currents).}
\question{Three identical $\qty{30}{\ohm}$ resistors in parallel are equivalent to:}
\opts{0}{\qty{90}{\ohm}}{0}{\qty{30}{\ohm}}{1}{\qty{10}{\ohm}}{0}{\qty{3.3}{\ohm}}
\why{$1/\Req = 3/30\ \mathrm{S}$, so $\Req = \qty{10}{\ohm}$.}

% ---------------------------------------------------------------------------
\exercise{Millman's theorem}{5 marks}
In the circuit below, the ground (bottom wire) is the reference of potentials. Using
Millman's theorem, determine the potential $V_A$ of node $A$. Deduce the current $i_3$ that
flows through $R_3$ towards the ground.
\begin{center}
\begin{circuitikz}[line width=0.6pt]
  \draw (0,0) to[battery1, invert, l=$E_1$] (0,1.4) to[R, l=$R_1$] (0,2.9);
  \draw (3,0) to[battery1, invert, l=$E_2$] (3,1.4) to[R, l=$R_2$] (3,2.9);
  \draw (6,2.9) to[R, l=$R_3$] (6,0);
  \draw (0,2.9) -- (3,2.9) -- (6,2.9);
  \draw (0,0) -- (3,0) -- (6,0);
  \node[circ] at (3,2.9) {}; \node[circ] at (3,0) {};
  \node[above] at (3,2.9) {$A$};
  \draw (3,0) node[cground, scale=1.4] {};
  \draw[-{Latex[length=1.8mm]}] (6,2.45) -- (6,2.15);
  \node[left] at (5.95,2.3) {$i_3$};
\end{circuitikz}
\end{center}
\data{$E_1 = \qty{12}{V}$ \hspace{6mm} $R_1 = \qty{2.0}{\ohm}$ \hspace{6mm}
      $E_2 = \qty{6.0}{V}$ \hspace{6mm} $R_2 = \qty{3.0}{\ohm}$ \hspace{6mm}
      $R_3 = \qty{6.0}{\ohm}$}
\answer{1,1,1,1,1}{6.8cm}{%
  The potentials at the far ends of the three branches are $E_1$, $E_2$ and $0$:\\[3pt]
  \hspace*{5mm}$\boxed{V_A = \dfrac{\dfrac{E_1}{R_1} + \dfrac{E_2}{R_2} + \dfrac{0}{R_3}}
     {\dfrac{1}{R_1} + \dfrac{1}{R_2} + \dfrac{1}{R_3}}}$\\[4pt]
  Numerator: $\dfrac{12}{2.0} + \dfrac{6.0}{3.0} + 0 = 6.0 + 2.0 = 8.0\ \mathrm{A}$ \qquad
  Denominator: $\dfrac{1}{2} + \dfrac{1}{3} + \dfrac{1}{6} = \dfrac{3 + 2 + 1}{6}
   = 1.0\ \mathrm{S}$\\[4pt]
  \hspace*{5mm}$\boxed{V_A = 8.0\ \mathrm{V}}$ \qquad
  $i_3 = \dfrac{V_A - 0}{R_3} = \dfrac{8.0}{6.0}\ \mathrm{A}$ \qquad
  $\boxed{i_3 = 1.3\ \mathrm{A}}$ \ ($4/3$~A)\\[4pt]
  Check (node law at $A$): $i_1 = \frac{E_1 - V_A}{R_1} = 2.0\ \mathrm{A}$ enters $A$,
  $\frac{E_2 - V_A}{R_2} = -0.67\ \mathrm{A}$, and $2.0 - 0.67 = 1.33\ \mathrm{A} = i_3$.
  \checkmark}
\begin{scheme}
\pt{1}{1 pt}{Millman's formula written literally, with the right potentials ($E_1$, $E_2$,
  $0$).}
\pt{2}{1 pt}{Numerator $= 8.0\ \mathrm{A}$.}
\pt{3}{1 pt}{Denominator $= 1.0\ \mathrm{S}$.}
\pt{4}{1 pt}{$V_A = 8.0\ \mathrm{V}$.}
\pt{5}{1 pt}{$i_3 = V_A/R_3 = 1.3\ \mathrm{A}$.}
\end{scheme}
\begin{tip}
Millman gives $V_A$ as a mean of the potentials weighted by the conductances. The result must
lie between the smallest and the largest potential ($0$ and $12\ \mathrm{V}$): a quick sanity
check.
\end{tip}

\finish
\end{document}
